\documentclass[11pt]{article}
\newcommand{\nombredocumento}{01 · Marco teórico}
% Preámbulo compartido por todos los documentos del proyecto (compilar con pdflatex).
\usepackage[utf8]{inputenc}
\usepackage[T1]{fontenc}
\usepackage{lmodern}
\usepackage{textcomp}
\usepackage[spanish,es-nodecimaldot,es-noshorthands,es-tabla]{babel}
\DeclareUnicodeCharacter{03A9}{\ensuremath{\Omega}}
\DeclareUnicodeCharacter{00B5}{\ensuremath{\mu}}
\DeclareUnicodeCharacter{03BC}{\ensuremath{\mu}}
\DeclareUnicodeCharacter{221A}{\ensuremath{\surd}}
\DeclareUnicodeCharacter{2192}{\ensuremath{\rightarrow}}
\DeclareUnicodeCharacter{2248}{\ensuremath{\approx}}
\DeclareUnicodeCharacter{2212}{\ensuremath{-}}
\DeclareUnicodeCharacter{2264}{\ensuremath{\le}}
\DeclareUnicodeCharacter{2265}{\ensuremath{\ge}}
\DeclareUnicodeCharacter{00D7}{\ensuremath{\times}}
\usepackage[a4paper,top=2.5cm,bottom=2.5cm,left=2.4cm,right=2.4cm,headheight=15pt]{geometry}
\setlength{\emergencystretch}{3em}
\usepackage{amsmath,amssymb,mathtools,bm}
\usepackage{siunitx}
\sisetup{output-decimal-marker={.},group-digits=false,per-mode=symbol,separate-uncertainty}
\usepackage{graphicx,float,booktabs,tabularx,longtable,array,multirow}
\usepackage[font=small,labelfont=bf]{caption}
\usepackage{enumitem}
\setlist{itemsep=2pt,topsep=4pt}
\usepackage[dvipsnames,table]{xcolor}
\usepackage[most]{tcolorbox}
\usepackage{tikz}
\usetikzlibrary{arrows.meta,positioning,shapes.geometric,calc,fit,backgrounds,decorations.markings}
\usepackage[european,straightvoltages]{circuitikz}
\usepackage{listings}
\usepackage{fancyhdr}
\usepackage{titlesec}
\usepackage{xurl}
\usepackage[colorlinks=true,linkcolor=Azul,urlcolor=Azul,citecolor=Azul]{hyperref}

\graphicspath{{../../simulacion/graficas/}{figuras/}}

% ---------------------------------------------------------------- colores
\definecolor{Azul}{HTML}{0D47A1}
\definecolor{Cian}{HTML}{00838F}
\definecolor{Naranja}{HTML}{E65100}
\definecolor{Rojo}{HTML}{C62828}
\definecolor{Verde}{HTML}{2E7D32}
\definecolor{Gris}{HTML}{455A64}
\definecolor{FondoCodigo}{HTML}{F5F7FA}

% ---------------------------------------------------------------- títulos
\titleformat{\section}{\Large\bfseries\color{Azul}}{\thesection}{0.8em}{}
\titleformat{\subsection}{\large\bfseries\color{Cian}}{\thesubsection}{0.8em}{}
\titleformat{\subsubsection}{\normalsize\bfseries\color{Gris}}{\thesubsubsection}{0.8em}{}

% ---------------------------------------------------------------- cabeceras
\pagestyle{fancy}
\fancyhf{}
\fancyhead[L]{\small\color{Gris}\nombredocumento}
\fancyhead[R]{\small\color{Gris}Generador de flujo axial con imanes NdFeB}
\fancyfoot[C]{\small\thepage}
\renewcommand{\headrulewidth}{0.4pt}

% ---------------------------------------------------------------- cajas
\newtcolorbox{nota}[1][Nota]{enhanced,breakable,colback=Azul!4,colframe=Azul!70,
  fonttitle=\bfseries,title={#1},left=6pt,right=6pt}
\newtcolorbox{advertencia}[1][Advertencia]{enhanced,breakable,colback=Rojo!4,colframe=Rojo!80,
  fonttitle=\bfseries,title={#1},left=6pt,right=6pt}
\newtcolorbox{resultado}[1][Resultado]{enhanced,breakable,colback=Verde!5,colframe=Verde!75,
  fonttitle=\bfseries,title={#1},left=6pt,right=6pt}
\newtcolorbox{paso}[1]{enhanced,breakable,colback=white,colframe=Cian!80,
  fonttitle=\bfseries,title={#1},left=6pt,right=6pt}

% ---------------------------------------------------------------- código
\lstdefinestyle{base}{
  basicstyle=\ttfamily\footnotesize,
  backgroundcolor=\color{FondoCodigo},
  frame=single,rulecolor=\color{Gris!40},
  breaklines=true,columns=fullflexible,keepspaces=true,
  keywordstyle=\color{Azul}\bfseries,
  commentstyle=\color{Verde}\itshape,
  stringstyle=\color{Naranja},
  showstringspaces=false,
  extendedchars=true,
  literate={á}{{\'a}}1 {é}{{\'e}}1 {í}{{\'i}}1 {ó}{{\'o}}1 {ú}{{\'u}}1 {ñ}{{\~n}}1
           {Á}{{\'A}}1 {É}{{\'E}}1 {Í}{{\'I}}1 {Ó}{{\'O}}1 {Ú}{{\'U}}1 {Ñ}{{\~N}}1
           {·}{{$\cdot$}}1 {→}{{$\rightarrow$}}1 {Ø}{{\O}}1 {Ω}{{$\Omega$}}1 {µ}{{$\mu$}}1
           {−}{{-}}1 {²}{{$^2$}}1 {√}{{$\surd$}}1 {≈}{{$\approx$}}1
}
\lstdefinestyle{python}{style=base,language=Python}
\lstdefinelanguage{JavaScript}{
  keywords={const,let,var,function,return,class,constructor,new,if,else,for,while,this,true,false,null,typeof,module,export,import},
  sensitive=true,comment=[l]{//},morecomment=[s]{/*}{*/},morestring=[b]',morestring=[b]",morestring=[b]`
}
\lstdefinestyle{js}{style=base,language=JavaScript}
\lstdefinestyle{consola}{style=base,keywordstyle=\color{black}}

% ---------------------------------------------------------------- utilidades
\newcommand{\vect}[1]{\mathbf{#1}}
\newcommand{\unitario}[1]{\mathbf{a}_{#1}}
\newcommand{\dif}{\mathrm{d}}
\DeclareUrlCommand\archivo{\urlstyle{tt}}
\makeatletter
\newcommand{\filasdesde}[1]{\@@input #1 }
\makeatother

\newcommand{\portada}[3]{%
  \begin{titlepage}
    \begin{tikzpicture}[remember picture,overlay]
      \fill[Azul] (current page.north west) rectangle ([yshift=-7cm]current page.north east);
      \node[anchor=north west,text=white,font=\sffamily\bfseries\Huge,text width=16cm,align=left]
        at ([xshift=2.4cm,yshift=-2.2cm]current page.north west) {#1};
      \node[anchor=north west,text=white!85!Azul,font=\sffamily\Large,text width=16cm,align=left]
        at ([xshift=2.4cm,yshift=-4.6cm]current page.north west) {#2};
    \end{tikzpicture}
    \vspace*{7cm}
    \begin{center}
      \begin{tikzpicture}[scale=0.9]
        \foreach \i in {0,...,7} {
          \pgfmathsetmacro{\ang}{45*\i}
          \ifodd\i \def\col{Azul!70}\def\pol{S}\else \def\col{Rojo!80}\def\pol{N}\fi
          \fill[\col] (\ang:2.2) circle (0.55);
          \node[white,font=\bfseries] at (\ang:2.2) {\pol};
        }
        \draw[Gris,line width=1pt] (0,0) circle (3.0);
        \fill[Gris!40] (0,0) circle (0.35);
      \end{tikzpicture}
    \end{center}
    \vfill
    {\large #3\par}
    \vspace{0.6cm}
    {\small\color{Gris}Basado en M.\,N.\,O. Sadiku, \emph{Elementos de Electromagnetismo}, 3.ª ed., Oxford University Press.\par
    Valores numéricos generados por el paquete \texttt{generador\_axial} (\texttt{uv run generador-axial reporte}).\par}
  \end{titlepage}
}

% Archivo generado automáticamente por `generador-axial reporte`. No editar a mano.
\newcommand{\ValGradoIman}{N38}  % Grado comercial del imán
\newcommand{\ValBr}{1.25}  % Remanencia [T]
\newcommand{\ValDiametroIman}{20}  % [mm]
\newcommand{\ValRadioIman}{10.0}  % [mm]
\newcommand{\ValEspesorIman}{5.0}  % [mm]
\newcommand{\ValAreaImanMM}{314.16}  % [mm^2]
\newcommand{\ValPolos}{8}
\newcommand{\ValParesPolos}{4}
\newcommand{\ValPasoPolarGrados}{45.0}  % [grados]
\newcommand{\ValRadioMontaje}{40.0}  % [mm]
\newcommand{\ValBobinas}{8}
\newcommand{\ValEspiras}{180}
\newcommand{\ValEspirasTotales}{1440}
\newcommand{\ValDiametroInteriorBobina}{12.0}  % [mm]
\newcommand{\ValDiametroExteriorBobina}{28.0}  % [mm]
\newcommand{\ValAlturaBobina}{6.0}  % [mm]
\newcommand{\ValEntrehierro}{2.5}  % [mm]
\newcommand{\ValEntrehierroMasAltura}{8.5}  % [mm]
\newcommand{\ValAcoplamiento}{0.85}
\newcommand{\ValCalibreAWG}{26}
\newcommand{\ValDiametroCobre}{0.404}  % [mm]
\newcommand{\ValDiametroEsmaltado}{0.455}  % [mm]
\newcommand{\ValAreaCobreMM}{0.1282}  % [mm^2]
\newcommand{\ValResistividadCobre}{1.72}  % [1e-8 ohm m]
\newcommand{\ValLongitudMediaEspira}{62.83}  % [mm]
\newcommand{\ValLongitudAlambreBobina}{11.31}  % [m]
\newcommand{\ValLongitudAlambreTotal}{90.48}  % [m]
\newcommand{\ValResistenciaBobina}{1.517}  % [ohm]
\newcommand{\ValResistenciaInterna}{12.14}  % [ohm]
\newcommand{\ValVentanaBobinaMM}{48.0}  % [mm^2]
\newcommand{\ValOcupacion}{0.78}
\newcommand{\ValRaizCaraIman}{11.1803}  % sqrt(L^2+R^2) [mm]
\newcommand{\ValTerminoCaraIman}{0.4472}  % L/sqrt(L^2+R^2)
\newcommand{\ValCampoCaraIman}{0.2795}  % [T]
\newcommand{\ValEntrehierroMasEspesor}{7.5}  % g+L [mm]
\newcommand{\ValRaizGL}{12.5000}  % sqrt((g+L)^2+R^2) [mm]
\newcommand{\ValRaizG}{10.3078}  % sqrt(g^2+R^2) [mm]
\newcommand{\ValTerminoUnoG}{0.6000}  % (g+L)/sqrt((g+L)^2+R^2)
\newcommand{\ValTerminoDosG}{0.2425}  % g/sqrt(g^2+R^2)
\newcommand{\ValCampoCaraBobina}{0.2234}  % [T]
\newcommand{\ValEntrehierroAlturaEspesor}{13.5}  % g+h+L [mm]
\newcommand{\ValRaizGHL}{16.8003}  % sqrt((g+h+L)^2+R^2) [mm]
\newcommand{\ValRaizGH}{13.1244}  % sqrt((g+h)^2+R^2) [mm]
\newcommand{\ValPrimitivaEnG}{2.1922}  % F(g) [mm]
\newcommand{\ValPrimitivaEnGH}{3.6759}  % F(g+h) [mm]
\newcommand{\ValDiferenciaPrimitivas}{1.4837}  % F(g+h)-F(g) [mm]
\newcommand{\ValCampoMedioBobina}{0.1545}  % [T]
\newcommand{\ValReduccionPorPromedio}{30.8}  % [%]
\newcommand{\ValFlujoMaximoMicro}{41.270}  % [uWb]
\newcommand{\ValAreaImanM}{3.1416}  % [1e-4 m^2]
\newcommand{\ValCaidaDiodo}{0.30}  % [V]
\newcommand{\ValCaidaPuente}{0.60}  % [V]
\newcommand{\ValCapacitanciaMicro}{220}  % [uF]
\newcommand{\ValTensionNominalCondensador}{50}  % [V]
\newcommand{\ValRpmMaxima}{1200}
\newcommand{\ValTensionDcMaxima}{29.27}  % [V]
\newcommand{\ValMargenCondensador}{41}  % [%]
\newcommand{\ValResistenciaLed}{2200}  % [ohm]
\newcommand{\ValPotenciaNominalResistenciaLed}{0.50}  % [W]
\newcommand{\ValCaidaLed}{2.0}  % [V]
\newcommand{\ValCorrienteLedRefMili}{2.21}  % [mA]
\newcommand{\ValCorrienteLedMaxMili}{12.4}  % [mA]
\newcommand{\ValCorrienteLedLimiteMili}{20}  % [mA]
\newcommand{\ValPotenciaResistenciaLedMax}{0.338}  % [W]
\newcommand{\ValRizadoRef}{251.4}  % [mV]
\newcommand{\ValRpmReferencia}{300}
\newcommand{\ValOmegaMecanicaRef}{31.416}  % [rad/s]
\newcommand{\ValFrecuenciaRef}{20.00}  % [Hz]
\newcommand{\ValPeriodoRefMs}{50.00}  % [ms]
\newcommand{\ValOmegaElectricaRef}{125.66}  % [rad/s]
\newcommand{\ValTensionPicoRef}{7.47}  % [V]
\newcommand{\ValTensionRmsRef}{5.28}  % [V]
\newcommand{\ValTensionDcRef}{6.87}  % [V]
\newcommand{\ValTensionPicoPorBobinaRef}{0.933}  % [V]
\newcommand{\ValRpmCarga}{400}
\newcommand{\ValOmegaMecanicaCarga}{41.89}  % [rad/s]
\newcommand{\ValTensionPicoCarga}{9.96}  % [V]
\newcommand{\ValTensionRmsCarga}{7.04}  % [V]
\newcommand{\ValCorrienteCarga}{0.290}  % [A]
\newcommand{\ValTensionEnCarga}{3.52}  % [V]
\newcommand{\ValPotenciaCarga}{1.021}  % [W]
\newcommand{\ValPotenciaPerdidas}{1.021}  % [W]
\newcommand{\ValPotenciaMecanica}{2.042}  % [W]
\newcommand{\ValParResistenteMili}{48.7}  % [mN m]
\newcommand{\ValEficiencia}{50}  % [%]
\newcommand{\ValBateriaTension}{12}  % [V]
\newcommand{\ValBateriaCapacidad}{7}  % [Ah]
\newcommand{\ValBateriaEnergia}{84}  % [Wh]
\newcommand{\ValBateriaResistencia}{0.05}  % [ohm]
\newcommand{\ValBateriaVacia}{11.8}  % E_oc con SoC = 0 [V]
\newcommand{\ValBateriaLlena}{12.8}  % E_oc con SoC = 1 [V]
\newcommand{\ValBateriaPendiente}{1.0}  % [V]
\newcommand{\ValSocInicial}{50}  % [%]
\newcommand{\ValSocCorte}{20}  % [%]
\newcommand{\ValSocReconexion}{25}  % [%]
\newcommand{\ValTensionOcRef}{12.30}  % [V]
\newcommand{\ValEnergiaUtilRef}{25.2}  % [Wh]
\newcommand{\ValControladorEficiencia}{90}  % [%]
\newcommand{\ValControladorArranque}{3.5}  % [V]
\newcommand{\ValControladorCorrienteMax}{3.0}  % [A]
\newcommand{\ValControladorAbsorcion}{14.4}  % [V]
\newcommand{\ValRpmMicrored}{600}
\newcommand{\ValFemRmsMicrored}{10.56}  % [V]
\newcommand{\ValTensionDcMicrored}{14.34}  % [V]
\newcommand{\ValPotenciaExtraibleMicrored}{2.297}  % [W]
\newcommand{\ValPotenciaHaciaBateria}{2.067}  % [W]
\newcommand{\ValCorrienteCargaMicroredMili}{168}  % [mA]
\newcommand{\ValPotenciaMecanicaMicrored}{4.59}  % [W]
\newcommand{\ValParMicroredMili}{73.1}  % [mN m]
\newcommand{\ValHorasRecarga}{20.8}  % [h]
\newcommand{\ValInversorTension}{110}  % [V]
\newcommand{\ValInversorPico}{155.6}  % [V]
\newcommand{\ValInversorTensionAlta}{220}  % [V]
\newcommand{\ValInversorPicoAlto}{311.1}  % [V]
\newcommand{\ValInversorFrecuencia}{60}  % [Hz]
\newcommand{\ValInversorNominal}{150}  % [W]
\newcommand{\ValInversorEficiencia}{90}  % [%]
\newcommand{\ValRelacionEspiras}{13.0}  % V_pico,AC / V_bat
\newcommand{\ValPotenciaCargaAc}{15}  % [W]
\newcommand{\ValCorrienteAcMili}{136}  % [mA]
\newcommand{\ValPotenciaEntradaInversor}{16.67}  % [W]
\newcommand{\ValCorrienteInversorBateria}{1.36}  % [A]
\newcommand{\ValPotenciaNetaBateria}{-14.60}  % [W]
\newcommand{\ValAutonomiaHoras}{1.73}  % [h]
\newcommand{\ValAutonomiaParadoHoras}{1.51}  % [h]
\newcommand{\ValRpmEquilibrio}{1704}
\newcommand{\ValPotenciaCargaReducida}{5}  % [W]
\newcommand{\ValRpmEquilibrioReducida}{984}
\newcommand{\ValInerciaRotor}{2.415}  % [1e-4 kg m^2]
\newcommand{\ValInductanciaBobinaMili}{0.528}  % [mH]
\newcommand{\ValInductanciaTotalMili}{4.22}  % [mH]
\newcommand{\ValFlujoMaximoReal}{0.827}  % phi_max normalizado
\newcommand{\ValPicoDerivadaFlujo}{0.746}
\newcommand{\ValRmsDerivadaFlujo}{0.559}
\newcommand{\ValFactorCresta}{1.335}
\newcommand{\ValFactorMedia}{0.706}
\newcommand{\ValRelacionFemRealIdeal}{0.930}
\newcommand{\ValPuntosRadiales}{32}
\newcommand{\ValMuestrasPerfil}{256}
\newcommand{\ValPasoIntegracionMili}{5}  % [ms]
\newcommand{\ValParCoulombMili}{3.0}  % [mN m]
\newcommand{\ValCoefViscosoMicro}{15.0}  % [uN m s]
\newcommand{\ValCoefVentilacionNano}{15.0}  % [nN m s^2]
\newcommand{\ValParMaximoMotriz}{0.30}  % [N m]
\newcommand{\ValVelocidadVacioMotriz}{1500}  % [RPM]
\newcommand{\ValParImpulso}{0.08}  % [N m]
\newcommand{\ValDuracionImpulso}{0.5}  % [s]
\newcommand{\ValCapacidadTermica}{100}  % [J/K]
\newcommand{\ValConveccionNatural}{8}  % [W/m^2K]
\newcommand{\ValConveccionForzada}{0.15}  % [W/m^2K por rad/s]
\newcommand{\ValCoefCobre}{0.393}  % [%/K]
\newcommand{\ValCoefRemanencia}{-0.12}  % [%/K]
\newcommand{\ValTemperaturaMaximaPla}{55}  % [C]
\newcommand{\ValBateriaRCero}{25}  % [mohm]
\newcommand{\ValBateriaRPolarizacion}{30}  % [mohm]
\newcommand{\ValBateriaTau}{30}  % [s]
\newcommand{\ValBateriaPeukert}{1.12}
\newcommand{\ValBateriaEtaCoulomb}{95}  % [%]
\newcommand{\ValBateriaEtaLlena}{60}  % [%]
\newcommand{\ValBateriaSocGas}{85}  % [%]
\newcommand{\ValBateriaAutodescarga}{3}  % [%/mes]
\newcommand{\ValBateriaSocFlotacion}{99.5}  % [%]
\newcommand{\ValBateriaFactorGas}{1500}
\newcommand{\ValBateriaPendienteGas}{50}
\newcommand{\ValControladorFlotacion}{13.8}  % [V]
\newcommand{\ValConsumoControladorMili}{8}  % [mA]
\newcommand{\ValInversorVacio}{3.6}  % [W]
\newcommand{\ValInversorCoefCuadratico}{5.81}  % [1e-4 1/W]
\newcommand{\ValInversorCorte}{10.8}  % [V]
\newcommand{\ValInversorReconexion}{12.3}  % [V]
\newcommand{\ValParReferencia}{0.07}  % [N m]
\newcommand{\ValRpmRegimenRef}{483}
\newcommand{\ValPotenciaEjeRef}{2.40}  % [W]
\newcommand{\ValPotenciaElectromagneticaRef}{2.21}  % [W]
\newcommand{\ValPerdidasCobreRef}{0.986}  % [W]
\newcommand{\ValPerdidasRozamientoRef}{0.192}  % [W]
\newcommand{\ValPerdidasFoucaultMili}{1.65}  % [mW]
\newcommand{\ValPotenciaBateriaRef}{0.96}  % [W]
\newcommand{\ValEficienciaRef}{40}  % [%]
\newcommand{\ValTensionDcRefReal}{3.50}  % [V]
\newcommand{\ValCorrienteGeneradorRef}{0.281}  % [A]
\newcommand{\ValTemperaturaRegimenRef}{27.8}  % [C]
\newcommand{\ValRpmConCargaRef}{483}
\newcommand{\ValPotenciaEntradaInversorRef}{18.7}  % [W]
\newcommand{\ValCorrienteNetaConCarga}{-1.48}  % [A]
\newcommand{\ValRpmParMaximo}{1002}
\newcommand{\ValPotenciaBateriaParMaximo}{4.35}  % [W]
\newcommand{\ValRizadoBajo}{18.8}  % [V] con 220 uF
\newcommand{\ValRizadoAlto}{0.88}  % [V] con 4700 uF
\newcommand{\ValRpmPicoImpulso}{505}
\newcommand{\ValTiempoParadaImpulso}{2.7}  % [s]
\newcommand{\ValHorasHastaAbsorcion}{4.0}  % [h]
\newcommand{\ValHorasHastaFlotacion}{8.2}  % [h]
\newcommand{\ValCorrienteFlotacionMili}{21}  % [mA]



\begin{document}

\portada{Marco teórico del generador de flujo axial}
        {Deducción paso a paso de las leyes electromagnéticas que gobiernan el prototipo}
        {Documento 1 de 5 · Fundamentos (Sadiku, Cap.~7, 8 y 9)}

\tableofcontents
\newpage

%=====================================================================
\section{Objetivo y alcance}
%=====================================================================

Este documento explica \emph{por qué} funciona el generador y \emph{de dónde salen} cada una de las
ecuaciones que usan el modelo de simulación (\archivo{src/generador_axial}), el visualizador 3D y la
guía de fabricación. Ninguna fórmula se presenta sin su deducción: se parte de la ley de Biot--Savart y de
la ley de Faraday y se llega, paso a paso, a la tensión que se medirá en los bornes del prototipo.

El generador estudiado tiene la siguiente topología:
\begin{itemize}
  \item \textbf{Rotor}: un disco con $P=\ValPolos$ imanes cilíndricos de NdFeB
        ($\varnothing\,\ValDiametroIman\times\ValEspesorIman$\,mm) montados a un radio de \ValRadioMontaje\,mm,
        con polaridad alternada N--S--N--S.
  \item \textbf{Estator}: un disco fijo con $N_b=\ValBobinas$ bobinas de núcleo de aire de $N=\ValEspiras$ espiras.
  \item \textbf{Entrehierro}: $g=\ValEntrehierro$\,mm entre la cara de los imanes y la cara de las bobinas.
\end{itemize}

\begin{nota}[Sobre la numeración de Sadiku]
Las secciones citadas corresponden a la 3.ª edición en español de \emph{Elementos de Electromagnetismo}:
Cap.~7 (Campos magnetostáticos: Biot--Savart, Ampère, Sec.~7.5 densidad de flujo),
Cap.~8 (Fuerzas, materiales y dispositivos magnéticos: Sec.~8.2 fuerzas, Sec.~8.5 imanación,
Sec.~8.6 clasificación de materiales) y Cap.~9 (Campos variables: Sec.~9.2 ley de Faraday,
Sec.~9.3 fem de transformador y de movimiento). En otras ediciones la numeración puede variar ligeramente.
\end{nota}

%=====================================================================
\section{Magnetostática: de Biot--Savart a la ley de Gauss magnética}
%=====================================================================

\subsection{Ley de Biot--Savart (Cap.~7)}
Un elemento de corriente $I\,\dif\vect{l}$ produce en un punto situado a un vector $\vect{R}$ de él la densidad
de flujo
\begin{equation}
  \dif\vect{B} = \frac{\mu_0}{4\pi}\,\frac{I\,\dif\vect{l}\times\unitario{R}}{R^2},
  \qquad \mu_0 = 4\pi\times10^{-7}\ \si{\henry\per\meter}.
  \label{eq:biot}
\end{equation}
Para una distribución superficial de corriente $\vect{K}$ (en \si{\ampere\per\meter}) basta sustituir
$I\,\dif\vect{l}\to\vect{K}\,\dif S$.

\subsection{Ley de Gauss para el campo magnético (Sec.~7.5)}
\begin{equation}
  \oint_S \vect{B}\cdot\dif\vect{S} = 0 \quad\Longleftrightarrow\quad \nabla\cdot\vect{B}=0 .
\end{equation}
\textbf{Paso a paso:} la forma diferencial se obtiene aplicando el teorema de la divergencia,
$\oint_S\vect{B}\cdot\dif\vect{S}=\int_V\nabla\cdot\vect{B}\,\dif v=0$, que debe cumplirse para \emph{cualquier}
volumen $V$; por tanto el integrando es nulo en todo punto.

\textbf{Consecuencia para el generador:} las líneas de $\vect{B}$ son cerradas. Todo el flujo que sale de la
cara Norte de un imán debe regresar a su cara Sur. En el rotor con polos alternados ese retorno se hace, en
buena parte, a través del imán vecino (de polaridad opuesta), lo que concentra el flujo en las bobinas que
están enfrente.

%=====================================================================
\section{Materiales magnéticos e imanes de neodimio}
%=====================================================================

\subsection{Imanación y corrientes ligadas (Sec.~8.5)}
En un material, los momentos magnéticos atómicos se describen con el vector imanación $\vect{M}$ (momento
dipolar por unidad de volumen, \si{\ampere\per\meter}). Se cumple
\begin{equation}
  \vect{B} = \mu_0\left(\vect{H}+\vect{M}\right).
\end{equation}
Sadiku demuestra que el efecto de $\vect{M}$ es idéntico al de dos distribuciones de corriente
\emph{ligada}:
\begin{equation}
  \vect{J}_b = \nabla\times\vect{M} \quad\text{(volumétrica)}, \qquad
  \vect{K}_b = \vect{M}\times\unitario{n} \quad\text{(superficial)}.
  \label{eq:ligadas}
\end{equation}

\subsection{Ferromagnéticos duros (Sec.~8.6)}
El NdFeB ($\mathrm{Nd_2Fe_{14}B}$) es un ferromagnético \emph{duro}: su ciclo de histéresis es muy ancho.
Las dos cifras que lo caracterizan son:
\begin{itemize}
  \item \textbf{Remanencia} $B_r$: densidad de flujo que conserva sin campo externo.
  \item \textbf{Coercitividad} $H_c$: campo inverso necesario para anular $B$ ($>\SI{850}{\kilo\ampere\per\meter}$).
        Es tan grande que las corrientes inducidas en las bobinas no pueden desmagnetizar el imán.
\end{itemize}

\begin{table}[H]
\centering
\caption{Grados comerciales típicos de NdFeB.}
\begin{tabular}{lccccc}
\toprule
Grado & N35 & N38 & N42 & N48 & N52\\
\midrule
$B_r$ típica [T] & 1.19 & 1.24 & 1.30 & 1.40 & 1.45\\
\bottomrule
\end{tabular}
\end{table}
El prototipo usa $B_r=\ValBr$\,T (grado \ValGradoIman, cuya remanencia comercial va de 1.22 a 1.26\,T).

%=====================================================================
\section{Campo axial de un imán cilíndrico: deducción completa}
%=====================================================================

\subsection{Paso 1: el imán es un solenoide equivalente}
Un imán cilíndrico de radio $R$ y espesor $L$ con imanación uniforme $\vect{M}=M\,\unitario{z}$ cumple:
\begin{enumerate}
  \item En el interior $\vect{M}$ es constante, así que $\vect{J}_b=\nabla\times\vect{M}=\vect{0}$.
  \item En las caras planas $\unitario{n}=\pm\unitario{z}$, de modo que $\vect{K}_b=M\unitario{z}\times(\pm\unitario{z})=\vect{0}$.
  \item En la cara lateral $\unitario{n}=\unitario{\rho}$, de modo que
        $\vect{K}_b=M\unitario{z}\times\unitario{\rho}=M\,\unitario{\phi}$.
\end{enumerate}
Es decir, \textbf{el imán equivale a una lámina cilíndrica de corriente azimutal $K=M$}: exactamente un
solenoide de longitud $L$ con $nI=M$ amperios-vuelta por metro.

\subsection{Paso 2: campo de una espira circular en su eje}
Tomemos una espira de radio $R$ en el plano $z'=0$ con corriente $I$, y un punto del eje a una altura $z$.
\begin{enumerate}
  \item La distancia de cualquier punto de la espira al punto de observación es $\sqrt{R^2+z^2}$, constante.
  \item Por simetría, las componentes radiales de $\dif\vect{B}$ se cancelan dos a dos; sólo sobrevive la axial.
  \item El ángulo entre $\dif\vect{l}$ y $\vect{R}$ es \ang{90}, y la proyección axial de $\dif\vect{B}$ introduce
        el factor $R/\sqrt{R^2+z^2}$:
        \[
          \dif B_z = \frac{\mu_0 I\,\dif l}{4\pi(R^2+z^2)}\cdot\frac{R}{\sqrt{R^2+z^2}} .
        \]
  \item Integrando $\dif l$ sobre la circunferencia ($\oint\dif l=2\pi R$):
        \begin{equation}
          B_z^{\text{espira}}(z) = \frac{\mu_0 I R^2}{2\left(R^2+z^2\right)^{3/2}} .
          \label{eq:espira}
        \end{equation}
\end{enumerate}

\subsection{Paso 3: superposición de espiras (solenoide finito)}
Situamos el imán entre $z'=-L$ y $z'=0$, de forma que su cara Norte está en $z=0$ y el espacio de trabajo es
$z>0$. Una rebanada de espesor $\dif z'$ transporta la corriente $K\,\dif z'$ y, por \eqref{eq:espira},
\[
  \dif B_z = \frac{\mu_0 K R^2\,\dif z'}{2\left(R^2+(z-z')^2\right)^{3/2}} .
\]
\begin{enumerate}
  \item Cambio de variable $u=z-z'$, $\dif u=-\dif z'$. Los límites cambian: $z'=-L\Rightarrow u=z+L$;
        $z'=0\Rightarrow u=z$.
  \item La integral queda
        $\displaystyle B_z=\frac{\mu_0 K R^2}{2}\int_{z}^{z+L}\frac{\dif u}{(R^2+u^2)^{3/2}}$.
  \item Primitiva conocida: $\displaystyle\int\frac{\dif u}{(R^2+u^2)^{3/2}}=\frac{u}{R^2\sqrt{R^2+u^2}}$
        (se comprueba derivando).
  \item Evaluando y simplificando $R^2$:
        \[
          B_z(z)=\frac{\mu_0 K}{2}\left[\frac{z+L}{\sqrt{(z+L)^2+R^2}}-\frac{z}{\sqrt{z^2+R^2}}\right].
        \]
  \item Para un imán permanente con permeabilidad de retroceso $\mu_r\approx1$ (NdFeB: 1.05), $\mu_0 M = B_r$.
        Con $K=M$:
\end{enumerate}
\begin{resultado}[Campo axial del imán]
\begin{equation}
  B_z(z)=\frac{B_r}{2}\left[\frac{z+L}{\sqrt{(z+L)^2+R^2}}-\frac{z}{\sqrt{z^2+R^2}}\right]
  \label{eq:campo}
\end{equation}
\end{resultado}

\subsection{Paso 4: comprobaciones de la fórmula}
\begin{itemize}
  \item \textbf{En la cara} ($z=0$): $B_z(0)=\dfrac{B_r}{2}\dfrac{L}{\sqrt{L^2+R^2}}$. Para el prototipo,
        $\ValCampoCaraIman$\,T.
  \item \textbf{Imán muy largo} ($L\to\infty$, $z=0$): $B_z\to B_r/2$, el valor conocido en el extremo de un solenoide
        semi-infinito.
  \item \textbf{Campo lejano} ($z\gg R,L$): desarrollando en serie se obtiene
        \[
          B_z\approx\frac{B_r R^2 L}{2\,d^3},
        \]
        con $d$ la distancia al centro del imán. Es el campo de un dipolo $m=M\pi R^2L$ en su eje,
        $B=\mu_0 m/(2\pi d^3)$. El paquete de pruebas verifica este límite.
\end{itemize}

\subsection{Paso 5: campo medio en el espesor de la bobina}
Las espiras de la bobina no están en un único plano: ocupan desde su cara ($z=g$) hasta su cara posterior
($z=g+h$), con $h=\ValAlturaBobina$\,mm. Lo correcto es promediar:
\[
  \bar B_z=\frac{1}{h}\int_g^{g+h}B_z(z)\,\dif z .
\]
\begin{enumerate}
  \item Primitivas de cada término de \eqref{eq:campo}:
        \[
          \int\frac{z+L}{\sqrt{(z+L)^2+R^2}}\,\dif z=\sqrt{(z+L)^2+R^2},\qquad
          \int\frac{z}{\sqrt{z^2+R^2}}\,\dif z=\sqrt{z^2+R^2}.
        \]
  \item Definimos $F(z)=\sqrt{(z+L)^2+R^2}-\sqrt{z^2+R^2}$.
\end{enumerate}
\begin{resultado}[Campo medio en la bobina]
\begin{equation}
  \bar B_z=\frac{B_r}{2h}\Big[F(g+h)-F(g)\Big]
  \label{eq:campomedio}
\end{equation}
\end{resultado}
\begin{nota}[Corrección respecto a la versión anterior del proyecto]
El modelo original evaluaba $B_z$ sólo en $z=g$ (\ValCampoCaraBobina\,T). El promedio exacto da
\ValCampoMedioBobina\,T: la versión anterior sobreestimaba el flujo, y por tanto la tensión, en un
\ValReduccionPorPromedio\,\%.
\end{nota}

%=====================================================================
\section{Flujo magnético y flujo concatenado}
%=====================================================================
El flujo a través de una superficie $S$ es $\Phi=\int_S\vect{B}\cdot\dif\vect{S}$ (\si{\weber}). Si una bobina
tiene $N$ espiras que abrazan prácticamente el mismo flujo, el \emph{flujo concatenado} es $\Psi=N\Phi$.

Para una espira enfrentada a un imán aproximamos
\begin{equation}
  \Phi_{max}=k\,\bar B_z\,\pi R^2,
  \label{eq:flujo}
\end{equation}
donde $\pi R^2$ es el área de la cara del imán y $k\in(0,1]$ es el \textbf{factor de acoplamiento}, que agrupa
lo que el modelo axial no recoge: la disminución de $B_z$ fuera del eje, el flujo de retorno que atraviesa las
espiras exteriores en sentido contrario y la dispersión en el entrehierro. Se usa $k=\ValAcoplamiento$.

%=====================================================================
\section{Ley de Faraday (Sec.~9.2)}
%=====================================================================
\begin{equation}
  \mathcal{E}=\oint_L\vect{E}\cdot\dif\vect{l}=-\frac{\dif\Psi}{\dif t}=-N\frac{\dif\Phi}{\dif t}.
  \label{eq:faraday}
\end{equation}
\textbf{Forma diferencial paso a paso.} Para un contorno fijo:
\begin{enumerate}
  \item $\oint_L\vect{E}\cdot\dif\vect{l}=-\dfrac{\dif}{\dif t}\displaystyle\int_S\vect{B}\cdot\dif\vect{S}
        =-\int_S\frac{\partial\vect{B}}{\partial t}\cdot\dif\vect{S}$ (el contorno no se mueve).
  \item Teorema de Stokes: $\oint_L\vect{E}\cdot\dif\vect{l}=\int_S(\nabla\times\vect{E})\cdot\dif\vect{S}$.
  \item Igualando para cualquier superficie $S$:
        \begin{equation}\nabla\times\vect{E}=-\frac{\partial\vect{B}}{\partial t}.\end{equation}
\end{enumerate}

%=====================================================================
\section{Fem de transformador y fem de movimiento (Sec.~9.3)}
%=====================================================================
Sadiku distingue tres casos:
\begin{align}
  \text{Espira fija, campo variable:}\quad &\mathcal{E}_{tr}=-\int_S\frac{\partial\vect{B}}{\partial t}\cdot\dif\vect{S},\\
  \text{Espira móvil, campo estático:}\quad &\mathcal{E}_{m}=\oint_L(\vect{u}\times\vect{B})\cdot\dif\vect{l},\\
  \text{Caso general:}\quad &\mathcal{E}=\mathcal{E}_{tr}+\mathcal{E}_{m}.
\end{align}
La fem de movimiento proviene de la fuerza de Lorentz $\vect{F}=q(\vect{E}+\vect{u}\times\vect{B})$: para una
carga que se mueve con el conductor, $\vect{u}\times\vect{B}$ actúa como un campo eléctrico efectivo.

\textbf{Aplicación al generador.}
\begin{itemize}
  \item \emph{Marco del estator}: las bobinas están quietas; en su posición el campo cambia al pasar polos
        alternos: es una fem de \textbf{transformador}.
  \item \emph{Marco del rotor}: el campo de los imanes es estático y son las bobinas las que se mueven con
        $\vect{u}=-\bm{\omega}_m\times\vect{r}$: es una fem de \textbf{movimiento}.
\end{itemize}
Ambas descripciones dan el mismo resultado: para $u\ll c$ las transformaciones de los campos entre marcos se
reducen a $\vect{E}'=\vect{E}+\vect{u}\times\vect{B}$, de modo que lo que en un marco aparece como
$\partial\vect{B}/\partial t$ en el otro aparece como $\vect{u}\times\vect{B}$.

%=====================================================================
\section{Modelo del generador}
%=====================================================================

\subsection{Paso polar y frecuencia eléctrica}
\begin{enumerate}
  \item Con $P$ polos repartidos uniformemente, el \textbf{paso polar} es $\theta_p=2\pi/P$
        (\ValPasoPolarGrados$^\circ$ para $P=\ValPolos$).
  \item Un ciclo eléctrico completo (N $\to$ S $\to$ N frente a una bobina) requiere girar dos pasos polares,
        $2\theta_p=4\pi/P$.
  \item En una vuelta ($2\pi$) hay por tanto $P/2$ ciclos eléctricos.
  \item Con $\text{RPM}/60$ vueltas por segundo:
\end{enumerate}
\begin{equation}
  f_e=\frac{P}{2}\cdot\frac{\text{RPM}}{60}=\frac{P\cdot\text{RPM}}{120},\qquad
  \omega_m=\frac{2\pi\,\text{RPM}}{60},\qquad \omega_e=2\pi f_e=\frac{P}{2}\,\omega_m .
  \label{eq:frecuencia}
\end{equation}

\subsection{Forma del flujo en una bobina}
El flujo que atraviesa una bobina es una función periódica del ángulo del rotor con periodo $2\theta_p$, y
cambia de signo cada $\theta_p$ (polo N $\to$ polo S). Su primer armónico de Fourier es
\begin{equation}
  \Phi(t)=\Phi_{max}\cos(\omega_e t),
\end{equation}
tomando $t=0$ cuando un imán N está centrado sobre la bobina. El modelo conserva sólo este armónico
fundamental, que es el que domina en topologías sin hierro con imanes circulares.

\subsection{Fem de una bobina}
Aplicando \eqref{eq:faraday}:
\begin{enumerate}
  \item $\dfrac{\dif\Phi}{\dif t}=-\Phi_{max}\,\omega_e\sin(\omega_e t)$ (regla de la cadena).
  \item $\mathcal{E}_1(t)=-N\dfrac{\dif\Phi}{\dif t}=N\,\omega_e\,\Phi_{max}\sin(\omega_e t)$.
\end{enumerate}

\subsection{Conexión en serie con polaridad alternada}
La bobina $j$ ($j=0,\dots,N_b-1$) está desplazada un ángulo $j\theta_p$, por lo que ve
$\Phi_j(t)=\Phi_{max}\cos(\omega_e t-j\pi)=(-1)^j\,\Phi_{max}\cos(\omega_e t)$: las bobinas vecinas tienen flujos
opuestos. Si todas se conectaran en el mismo sentido, sus fem se cancelarían por pares. Conectando la bobina $j$
con signo $s_j=(-1)^j$ (final con final, inicio con inicio):
\[
  \mathcal{E}_{total}=\sum_{j=0}^{N_b-1}s_j\,\mathcal{E}_j=\sum_{j}(-1)^j(-1)^j\,N\omega_e\Phi_{max}\sin(\omega_e t)
  =N_b N\,\omega_e\Phi_{max}\sin(\omega_e t).
\]
\begin{resultado}[Fem del generador]
\begin{equation}
  \mathcal{E}(t)=V_{pico}\sin(\omega_e t),\qquad V_{pico}=N_{total}\,\omega_e\,\Phi_{max},\qquad N_{total}=N_bN=\ValEspirasTotales .
  \label{eq:fem}
\end{equation}
\end{resultado}

\subsection{Valor eficaz}
Por definición $V_{RMS}^2=\frac1T\int_0^T\mathcal{E}^2\dif t$. Usando $\sin^2x=\frac{1-\cos2x}{2}$ y que el coseno
promedia cero en un periodo:
\begin{equation}
  V_{RMS}^2=\frac{V_{pico}^2}{T}\int_0^T\frac{1-\cos(2\omega_e t)}{2}\dif t=\frac{V_{pico}^2}{2}
  \;\Rightarrow\; V_{RMS}=\frac{V_{pico}}{\sqrt2}.
\end{equation}

%=====================================================================
\section{Circuito equivalente, potencia y ley de Lenz}
%=====================================================================

\subsection{Resistencia interna}
Para una bobina con diámetros interior $D_i$ y exterior $D_e$, la espira media tiene diámetro $(D_i+D_e)/2$:
\begin{equation}
  \ell_m=\pi\frac{D_i+D_e}{2},\qquad R_{int}=\rho_{Cu}\frac{N_{total}\,\ell_m}{A_{Cu}},\qquad A_{Cu}=\frac{\pi d^2}{4}.
\end{equation}
\begin{nota}[Corrección]
La versión anterior usaba un radio medio de 12\,mm; con $D_i=12$\,mm y $D_e=28$\,mm el radio medio real es
10\,mm. Con el valor corregido, $R_{int}=\ValResistenciaInterna\,\Omega$.
\end{nota}
Al ser bobinas sin núcleo su inductancia es pequeña y $Z_{int}=R_{int}+j\omega_eL_{int}\approx R_{int}$.

\subsection{Teorema de máxima transferencia de potencia}
Con una carga resistiva $R_L$ (circuito de Thévenin, valores eficaces):
\begin{enumerate}
  \item $I=\dfrac{V_{RMS}}{R_{int}+R_L}$, \quad $P_L=I^2R_L=\dfrac{V_{RMS}^2R_L}{(R_{int}+R_L)^2}$.
  \item Derivando respecto a $R_L$:
        $\dfrac{\dif P_L}{\dif R_L}=V_{RMS}^2\dfrac{(R_{int}+R_L)-2R_L}{(R_{int}+R_L)^3}=V_{RMS}^2\dfrac{R_{int}-R_L}{(R_{int}+R_L)^3}$.
  \item Se anula para $R_L=R_{int}$ (es un máximo: la derivada pasa de positiva a negativa).
  \item $P_{max}=\dfrac{V_{RMS}^2}{4R_{int}}$, con rendimiento $\eta=\dfrac{R_L}{R_{int}+R_L}$, que vale \SI{50}{\percent}.
\end{enumerate}

\subsection{Ley de Lenz y par resistente}
El signo menos de \eqref{eq:faraday} indica que la corriente inducida crea un flujo que se opone a la variación
que la produjo. Mecánicamente:
\begin{enumerate}
  \item Cada tramo de conductor con corriente sufre $\dif\vect{F}=I\,\dif\vect{l}\times\vect{B}$ (Sec.~8.2).
  \item La reacción sobre los imanes produce un par $\bm\tau=\vect{r}\times\vect{F}$ opuesto a $\bm\omega_m$.
  \item Conservación de la energía: la potencia mecánica entregada al eje iguala la potencia eléctrica total
        \begin{equation}
          P_{mec}=\tau\,\omega_m=I^2(R_{int}+R_L)=P_L+P_{Cu}\quad\Rightarrow\quad\tau=\frac{P_L+P_{Cu}}{\omega_m}.
        \end{equation}
\end{enumerate}
No hay energía gratuita: cuanto más corriente se extrae, más cuesta girar el rotor.

%=====================================================================
\section{Rectificación y filtrado}
%=====================================================================
\begin{enumerate}
  \item En el puente de Graetz conducen dos diodos en serie en cada semiciclo; la tensión de pico disponible en la
        salida es $V_{pico}-2V_D$.
  \item El condensador $C$ se carga al valor de pico y, entre picos, alimenta la carga. Con rectificación de onda
        completa los picos se repiten cada $1/(2f_e)$.
  \item Si la carga extrae una corriente casi constante $I_{DC}$, el condensador pierde
        $\Delta Q=I_{DC}/(2f_e)$ y su tensión cae $\Delta V=\Delta Q/C$:
        \begin{equation}
          V_{DC}\approx V_{pico}-2V_D,\qquad \Delta V_{pp}\approx\frac{I_{DC}}{2f_eC}.
        \end{equation}
\end{enumerate}
Con diodos 1N4007 $V_D\approx0.7$\,V; con Schottky 1N5819 $V_D\approx\ValCaidaDiodo$\,V (valor del modelo).

%=====================================================================
\section{¿Por qué un estator sin núcleo de hierro?}
%=====================================================================
\begin{itemize}
  \item \textbf{Con hierro}: $\mu_r\sim10^3$ aumenta el flujo, pero aparecen pérdidas por histéresis
        ($\propto f\,B^{n}$), por corrientes parásitas ($\propto f^2B^2$) y un fuerte \emph{par de reluctancia}
        (\emph{cogging}): el rotor ``se engancha'' en posiciones preferentes.
  \item \textbf{Sin hierro (coreless)}: $\mu_r\approx1$, cero cogging, sin pérdidas en el núcleo, sin saturación.
        El sistema es lineal y la ley de Faraday se verifica de forma limpia: ideal para un laboratorio.
\end{itemize}

%=====================================================================
\section{Micro-red: batería, controlador de carga e inversor}
%=====================================================================
El generador alimenta una micro-red aislada: controlador DC-DC $\rightarrow$ batería de
\ValBateriaTension\,V / \ValBateriaCapacidad\,Ah $\rightarrow$ inversor de \ValInversorTension\,V o
\ValInversorTensionAlta\,V a \ValInversorFrecuencia\,Hz $\rightarrow$ bombillo de \ValPotenciaCargaAc\,W.
La física electromagnética termina en la tensión del generador; el resto es teoría de circuitos y de
electrónica de potencia, salvo el transformador del inversor, que vuelve a ser la ley de Faraday.

\subsection{Modelo de la batería}
\begin{enumerate}
  \item Se modela como una fuente $E_{oc}$ que depende del estado de carga (SoC) en serie con una resistencia
        $R_{bat}=\ValBateriaResistencia\,\Omega$:
        \begin{equation}
          E_{oc}(\text{SoC})=\ValBateriaVacia+\ValBateriaPendiente\cdot\text{SoC}\ \text{V},\qquad
          V_{bat}=E_{oc}+I_{neta}R_{bat}.
        \end{equation}
  \item El SoC se integra con la corriente neta (conteo de Coulomb), con $Q=C\cdot\SI{3600}{\second}$:
        \begin{equation}
          \frac{\dif\,\text{SoC}}{\dif t}=\frac{I_{carga}-I_{descarga}}{Q}.
        \end{equation}
  \item Protección por baja tensión: el inversor se desconecta al \ValSocCorte\,\% y se reconecta al
        \ValSocReconexion\,\%. La histéresis evita que la salida oscile en el umbral.
\end{enumerate}

\subsection{Controlador de carga con seguimiento del punto de máxima potencia}
El generador es una fuente de Thévenin con $R_{int}=\ValResistenciaInterna\,\Omega$. Un controlador DC-DC que
emula una carga $R_L=R_{int}$ extrae, por el teorema de máxima transferencia,
\begin{equation}
  P_{max}=\frac{V_{RMS}^2}{4R_{int}},\qquad
  I_{carga}=\min\!\left(\frac{\eta_c\,P_{max}}{E_{oc}(\text{SoC})},\ I_{max},\ \frac{V_{abs}-E_{oc}}{R_{bat}}\right),
\end{equation}
con $\eta_c=\ValControladorEficiencia$\,\%, $I_{max}=\ValControladorCorrienteMax$\,A y
$V_{abs}=\ValControladorAbsorcion$\,V. El controlador no arranca si $V_{DC}<\ValControladorArranque$\,V.
\begin{nota}[Corrección]
La versión modificada del proyecto calculaba la potencia con $V_{DC}^2/(4R_{int})$. Como
$V_{DC}\approx V_{pico}=\sqrt2\,V_{RMS}$, eso duplicaba la potencia disponible. La forma correcta usa el valor
eficaz de la fem.
\end{nota}

\subsection{Inversor y transformador elevador (Sec.~8.8 y 9.3)}
\begin{enumerate}
  \item Un puente de MOSFET conmuta la batería con modulación senoidal (SPWM); el filtro LC deja sólo el
        armónico fundamental de \ValInversorFrecuencia\,Hz.
  \item El transformador eleva la tensión por la fem de transformador: ambos devanados comparten el flujo
        $\Phi(t)$, y por Faraday $v_1=N_1\,\dif\Phi/\dif t$ y $v_2=N_2\,\dif\Phi/\dif t$. Por tanto
        \begin{equation}
          \frac{V_2}{V_1}=\frac{N_2}{N_1}\approx\frac{\ValInversorPico\ \text{V}}{\ValBateriaTension\ \text{V}}
          \approx\ValRelacionEspiras .
        \end{equation}
  \item Balance de potencia con eficiencia $\eta_{inv}=\ValInversorEficiencia$\,\%:
        \begin{equation}
          P_{in}=\frac{P_{AC}}{\eta_{inv}},\qquad I_{descarga}=\frac{P_{AC}}{\eta_{inv}\,E_{oc}}.
        \end{equation}
\end{enumerate}

\subsection{Balance energético y velocidad de equilibrio}
La potencia neta que recibe la batería es
\begin{equation}
  P_{neta}=\eta_c\frac{V_{RMS}^2}{4R_{int}}-\frac{P_{AC}}{\eta_{inv}}.
\end{equation}
Como $V_{RMS}\propto\text{RPM}$, la velocidad que equilibra el consumo ($P_{neta}=0$) es
\begin{equation}
  \text{RPM}_{eq}=\text{RPM}_{ref}\,\frac{\sqrt{4R_{int}P_{AC}/(\eta_c\eta_{inv})}}{V_{RMS}(\text{RPM}_{ref})}.
\end{equation}
\begin{resultado}
Para el bombillo de \ValPotenciaCargaAc\,W se necesitan $\approx\ValRpmEquilibrio$\,RPM, por encima del rango
del prototipo (1200\,RPM). Por tanto el generador \textbf{recarga} la batería y es la batería la que
\textbf{alimenta} el inversor. Una carga de \ValPotenciaCargaReducida\,W se equilibra a
$\approx\ValRpmEquilibrioReducida$\,RPM.
\end{resultado}

%=====================================================================
\section{Hipótesis y limitaciones del modelo}
%=====================================================================
\begin{enumerate}
  \item Se usa el campo \emph{en el eje} del imán; fuera del eje $B_z$ es menor (lo absorbe $k$).
  \item Se conserva sólo el armónico fundamental del flujo.
  \item Se desprecian la inductancia de las bobinas y la variación de $\rho_{Cu}$ con la temperatura
        ($\approx\SI{0.39}{\percent\per\kelvin}$).
  \item El rotor es de una sola cara y sin placa de retorno de hierro; un rotor doble (imanes a ambos lados del
        estator) aproximadamente duplicaría el flujo.
  \item $k$ debe calibrarse con la medición: $k_{real}=k\cdot V_{\text{medido}}/V_{\text{teórico}}$.
\end{enumerate}

%=====================================================================
\section{Resumen de fórmulas}
%=====================================================================
\begin{table}[H]
\centering
\renewcommand{\arraystretch}{1.6}
\begin{tabular}{p{5.2cm}p{7.6cm}c}
\toprule
Magnitud & Expresión & Unidad\\
\midrule
Velocidad angular mecánica & $\omega_m=2\pi\,\text{RPM}/60$ & rad/s\\
Frecuencia eléctrica & $f_e=P\cdot\text{RPM}/120$ & Hz\\
Campo axial del imán & $B_z=\frac{B_r}{2}\Big[\frac{z+L}{\sqrt{(z+L)^2+R^2}}-\frac{z}{\sqrt{z^2+R^2}}\Big]$ & T\\
Campo medio en la bobina & $\bar B_z=\frac{B_r}{2h}\big[F(g+h)-F(g)\big]$ & T\\
Flujo máximo por espira & $\Phi_{max}=k\,\bar B_z\,\pi R^2$ & Wb\\
Fem pico total & $V_{pico}=N_{total}\,\omega_e\,\Phi_{max}$ & V\\
Fem eficaz & $V_{RMS}=V_{pico}/\sqrt2$ & V\\
Tensión rectificada & $V_{DC}\approx V_{pico}-2V_D$ & V\\
Rizado & $\Delta V_{pp}\approx I_{DC}/(2f_eC)$ & V\\
Resistencia interna & $R_{int}=\rho_{Cu}N_{total}\pi(D_i+D_e)/(2A_{Cu})$ & $\Omega$\\
Potencia en la carga & $P_L=V_{RMS}^2R_L/(R_{int}+R_L)^2$ & W\\
Par resistente (Lenz) & $\tau=(P_L+P_{Cu})/\omega_m$ & N·m\\
Tensión de la batería & $E_{oc}=\ValBateriaVacia+\ValBateriaPendiente\cdot\text{SoC}$ & V\\
Corriente de carga (MPPT) & $I_{carga}=\eta_cV_{RMS}^2/(4R_{int}E_{oc})$ & A\\
Relación del transformador & $V_2/V_1=N_2/N_1$ & --\\
Corriente del inversor & $I_{descarga}=P_{AC}/(\eta_{inv}E_{oc})$ & A\\
\bottomrule
\end{tabular}
\end{table}

Los valores numéricos de cada una de estas magnitudes para el prototipo se calculan paso a paso en el
\textbf{Documento 2: Memoria de cálculo}.

\end{document}
